%% kapa: exam=1 round=32 year=2021 period=1 form=A subject=civil_law subject_ko=민법 seq=38 number=38
%% source: Q-Net 감정평가사 기출문제 (cst00309) artlSeq=5208397 file=extracted/1008/5208397/A형/2021년도 제32회 감정평가사 1차 1교시 A형.pdf
%% conversion: kordoc 4.24.51 -> tools/gen1.py (auto); review=proofread-sonnet; flags=none
\begin{question}{38}
저당권에 관한 설명으로 옳지 않은 것은? (다툼이 있으면 판례에 따름)
\begin{choices}
  \item 저당부동산에 대한 압류 후에는 저당권설정자의 저당부동산에 관한 차임채권에도 저당권의 효력이 미친다.
  \item 저당목적물의 변형물에 대하여 이미 제3자가 압류하였더라도 저당권자가 스스로 이를 압류하지 않으면 물상대위권을 행사할 수 없다.
  \item 저당권은 그 담보한 채권과 분리하여 타인에게 양도하거나 다른 채권의 담보로 하지 못한다.
  \item 저당권의 효력은 원칙적으로 저당부동산에 부합된 물건에 미친다.
  \item 저당부동산에 대하여 지상권을 취득한 제3자는 저당권자에게 그 부동산으로 담보된 채권을 변제하고 저당권의 소멸을 청구할 수 있다.
\end{choices}
\end{question}
\answer{2}
